\section{Evaluating Geographic Transfer Failure}
\label{sec:evaluation}

\subsection{Framework}
\label{sec:framework}

\begin{figure*}[t]
  \centering
  \includegraphics[height=4cm, width=\linewidth]{figure2.png}
  \caption{\textbf{Failure mode and recovery under geographic transfer.}
Rows: three architectures on tiles from their leave-one-city-out (LOCO) rotations, together spanning all three held-out cities. Columns: input, ground-truth mask (white $=$ shadow), and predicted shadow probability under Upper bound (within-city training), LOCO, and SIB (\S\ref{sec:sib}); shared colormap $[0, 1]$. $\bar{p}_S$: mean predicted shadow probability inside ground-truth shadow. \emph{What to look for:} LOCO collapses $\bar{p}_S$ while background stays dark — failure is washed-out predictions over true shadow regions, not hallucinated shadow elsewhere. SIB recovers foreground confidence on the same regions across all three architectures with varying magnitude.}
  \label{fig:qualitative}
\end{figure*}




\begin{figure*}[t]
  \centering
  \includegraphics[width=\linewidth]{fig3_2.jpeg}
  \caption{\textbf{Diagnostic framework.}
\textbf{(a)} Risk-coverage curves on foreground pixels (MAMNet, Phoenix held out): LOCO accumulates $\Delta$AURC$=+0.26$ over the upper bound; FDA narrows it. Bands: 95\% bootstrap CIs. \emph{Inset:} foreground gap replicates on INRIA building footprints under matched LOCO (C/M/P: Chicago/Miami/Phoenix; A/C/V: Austin/Chicago/Vienna).
\textbf{(b)} Per-method AURC gap across $9$ cells ($3$ architectures $\times$ $3$ cities): the gap is class-asymmetric, concentrated on foreground (coral) with a pooled background gap (black) $24\times$ smaller over the $63$ cells of the $7$ methods.
\textbf{(c)} Per-image recovery ratio $R$ under decoder (Exp~A) vs.\ encoder (Exp~B) retraining: the DINOv3$\leftrightarrow$OGLANet sign reversal shows the failure locus is architecture-dependent.
LOCO: leave-one-city-out; AURC: area under risk-coverage.}
  \label{fig:diagnostics}
\end{figure*}

Table~\ref{tab:loco-main} showed that no adaptation method is uniformly safe across architectures, suggesting the failure is not a single mechanism a single correction can address. Yet the gap is consistent enough within each architecture that we can ask \emph{where} in the model it lives. We answer this with three controlled interventions on the LOCO models from Table~\ref{tab:loco-main}, ordered by what a practitioner would try first. Section~\ref{sec:feature-space} tests the cheapest possible fix: aligning pixel statistics across cities, ruling in or out a whole class of preprocessing-based recoveries. Section~\ref{sec:sp-gap} then asks whether the model's outputs are still usable under a confidence threshold, since deployed systems consume thresholded predictions rather than raw probabilities. Section~\ref{sec:architecture-locus} localizes \emph{which component} of each architecture carries the failure. Figure~\ref{fig:qualitative} grounds what we are explaining: under LOCO, models keep spatial localization but lose foreground-class confidence: sharp upper-bound predictions become washed-out heatmaps over true shadow regions, not hallucinated shadow elsewhere.

\subsection{Failure operates in feature space, not pixel space}
\label{sec:feature-space}

\textbf{Test design.} If the LOCO failure were a pixel-statistics shortcut, aligning the test image's pixel histogram to the training-city distribution should at least partially recover the gap. We test this directly. For each (architecture, holdout city) pair under vanilla LOCO, we compute the aggregate intensity histogram (the mean of the three channels) of the source-city training images. We map each test image's intensity CDF onto this histogram, restore each pixel's ratio between its channels and its intensity, and run inference with the trained model unchanged. 

\textbf{Result.} Histogram matching makes things worse, uniformly. Across the $9$ vanilla LOCO cells, mean mIoU drops by $0.45$ relative to LOCO itself (paired bootstrap $95\%$ CI $[-0.523, -0.380]$, $p < 0.001$); not a single cell improves. To put the magnitude on a scale, we use the per-image recovery ratio $R_i = \frac{\text{IoU}_i^{\text{match}} - \text{IoU}_i^{\text{LOCO}}}{\text{IoU}_i^{\text{upper}} - \text{IoU}_i^{\text{LOCO}}}$ where $R = 1$ means full recovery; $R = 0$ means no change from LOCO; $R < 0$ means the intervention is worse than doing nothing. Per-architecture medians fall in $[-2.72, -1.55]$, with all $95\%$ CIs excluding zero: at the median, matched images lie $2.5$ to $3.7$ times as far from the upper bound as LOCO itself. The same rule improves none of $27$ vanilla LOCO models trained on $225$, $450$ or $900$ images (three architectures, high resolution; Appendix~\ref{app:ladder}).

\textbf{Training-time appearance augmentation.} We also retrain both the LOCO model and its upper bound with appearance-only augmentation (colour, illumination, contrast, blur, noise; geometry preserved) on the $9$ high-resolution (architecture, held-out city) cells. Retraining with a different random seed alone moves a cell's gap: across three vanilla seeds, the standard deviation of the gap ranges from $0.24$ to $4.26$ mIoU between cells. Before training the augmented models, we therefore set each cell's threshold at twice this standard deviation and count a gap as changed only if augmentation moves it by more. Six of nine gaps do not change. The mean gap moves from $-6.04$ to $-4.11$ mIoU, with OGLANet/Phoenix and DINOv3/Phoenix narrowing and DINOv3/Miami widening. Phoenix also has the largest seed variance of the three cities, and three cities cannot separate a Phoenix-specific appearance effect from that variance (Appendix~\ref{app:aug-stress}).

\textbf{Implication.} Geographic transfer failure is not a pixel-level shortcut. The features the LOCO model relies on encode the joint distribution of intensity, texture, and spatial pattern in source-city imagery. Aligning the marginal pixel histogram alone breaks this coupling without replacing it with anything useful. This rules out a whole class of cheap fixes a practitioner might try first (like gamma correction, contrast normalization, single-channel histogram standardization) and tells us any successful intervention has to operate at the feature level. 

\subsection{Class-asymmetric selective-prediction failure}
\label{sec:sp-gap}

\textbf{Test design.} A practitioner deploying segmentation under domain shift consumes confidence-thresholded predictions, not raw outputs. The relevant question is therefore not the average error but how the threshold-conditioned error rate degrades under shift. The selective-prediction (SP) framework~\cite{geifman2017selective, traub2024overcoming, cattelan2023fix} measures threshold-conditioned error: per image, sort pixels by confidence and compute the cumulative error rate at each coverage level $c \in \{0.10, 0.15, \ldots, 1.00\}$. The Area Under the Risk-Coverage curve (AURC) integrates this curve. The \emph{SP gap} = AURC$_{\text{LOCO}}$ $-$ AURC$_{\text{upper}}$ measures how much threshold-conditioned error rises under transfer; a gap of zero means selective deployment recovers within-city quality. We compute the gap on three pixel populations: the foreground class (GT-shadow, confidence $= p$), the background class (GT-background, confidence $= 1{-}p$), and all pixels (confidence $= \max(p, 1{-}p)$). All numbers in this subsection come from the $63$ LOCO cells (architecture $\times$ method $\times$ holdout city, for high-resolution); paired image-level bootstrap ($B = 1{,}000$) gives $95\%$ confidence intervals.

\textbf{The gap concentrates on the foreground class.} On shadow pixels the gap averages $+0.136$ across the $63$ cells. On background pixels it averages $+0.006$. The asymmetry is consistent across cells: the foreground gap exceeded the background gap in $55$ of the $63$ cells (sign-test $p = 4.9 \times 10^{-10}$). Two further measurements describe how the model's confidence shifts. At the logit level, true-positive logits drop on average by $0.85$ nats (the model becomes much less confident on its foreground predictions), while logits on false-positive predictions barely move ($-0.15$ nats); the drop is larger on the true-positive side in $55$ of $63$ cells. Calibration error rises asymmetrically in the same direction (Appendix~\ref{app:sp-recovery}). The LOCO failure is therefore not uniform confidence flattening: the model selectively loses commitment on the class it is trying to detect, while leaving its background commitment intact.

\textbf{Cross-task scope test on INRIA building footprints.} A natural objection is that these patterns are properties of shadow appearance, not geographic transfer. We replicate the protocol on INRIA building extraction~\cite{inria}, a dense-prediction task whose foreground is shape-defined rather than illumination-defined (MAMNet, holdout cities Austin/Chicago/Vienna, matched training recipe). The foreground SP gap and the foreground-vs-background asymmetry both transfer in direction across tasks (sign test at the strongest result the sample size permits, $p=0.016$ each; Cohen's $d>1.1$), confirming the SP-gap signature is a property of LOCO across tasks rather than of shadows. Figure~\ref{fig:diagnostics}(a) inset shows the gap distributions side by side.

\textbf{A recoverable threshold exists.} The selective-prediction framing is only useful if there is some confidence threshold at which LOCO predictions match within-city quality. We find one in $59$ of $63$ cells; the four exceptions are all on Phoenix, the most distant holdout, under the most aggressive adaptation methods. This makes coverage thresholding, rather than post-hoc temperature scaling, the natural label-free recovery recipe under geographic shift, and Section~\ref{sec:sib} builds on it directly. Method-ranking stability across architectures, the post-hoc temperature-scaling alternative, and per-cell threshold values are in Appendix~\ref{app:sp-recovery}.

\textbf{Implication.} The SP gap is large, foreground-concentrated, hidden by $0/1$ evaluation, recoverable by coverage thresholding, and consistent across architectures and across a second dense-prediction task. Section~\ref{sec:architecture-locus} asks where in the architecture the failure that produces this gap actually lives.

\subsection{Architecture-dependent failure locus}
\label{sec:architecture-locus}

Sections~\ref{sec:feature-space} and~\ref{sec:sp-gap} established that the failure is in feature space and that confidence collapses asymmetrically. Neither result tells us \emph{which part} of a given model produces the collapse: the encoder, the decoder, or both. This matters because uniform feature-level corrections (ISW whitens, MRFP$^+$ randomizes) act on a fixed locus, and Table~\ref{tab:loco-main} already showed they fail catastrophically on some architectures while helping others. If different architectures carry the failure in different components, that asymmetry is the explanation.

\textbf{Test design.} We localise the failure by targeted retraining. \textbf{Experiment A} freezes the LOCO-trained encoder, reinitializes the decoder, and retrains \emph{only} the decoder on $25\%$ of the holdout-city training set. \textbf{Experiment B} does the inverse: freezes the decoder, resets the encoder to its pretrained initialization, and retrains \emph{only} the encoder on the same subset. If a component carries the failure, retraining it on a small target sample should recover most of the within-city gap; retraining the other component should not. Recovery is measured by the per-image ratio $R$ defined earlier (training details in Appendix~\ref{app:locus-retraining}).

\textbf{Mirror-image recovery pattern.} Decoder retraining (Experiment A) recovers DINOv3 ($R = 0.69$) and MAMNet ($R = 0.67$) but barely moves OGLANet ($R = 0.12$, $[0.05, 0.21]$ - CI excludes meaningful recovery). Encoder retraining (Experiment B) shows the opposite pattern: it makes DINOv3 \textit{worse} than LOCO ($R = -0.13$), while recovering both CNNs ($R = 0.88$ MAMNet, $R = 0.86$ OGLANet). All CIs in Appendix~\ref{app:locus-retraining}. The DINOv3-vs-OGLANet sign reversal across the two experiments shows decoder works for one, encoder works for the other, and each architecture is harmed by the other's fix. This is the central evidence that the failure component is architecture-dependent. Figure~\ref{fig:diagnostics}(c) shows the pattern.

\textbf{Frozen encoder.} If the gap came only from fine-tuning the encoder on two source cities, freezing the encoder would remove it. With the DINOv3 encoder frozen and only the decoder trained, the upper bound still exceeds LOCO by $7.10$, $2.67$ and $2.75$ mIoU on Chicago, Miami and Phoenix (each significant after Holm correction over the three cities; Appendix~\ref{app:frozen}). Both models share the frozen encoder, so this gap arises in the decoder, consistent with the locus Experiment~A measures for DINOv3.

\textbf{Pretraining and fine-tuning regime.} DINOv3 also differs from the CNNs in self-supervised pretraining and a low fine-tuning learning rate. We test both on MAMNet with a ResNet-50 encoder by training four versions: ImageNet-supervised or self-supervised DINO pretraining, each fine-tuned at learning rate $10^{-4}$ or $10^{-5}$ (three held-out cities, three seeds each). Encoder retraining recovers all four versions (median $R$ $0.80$ to $0.87$, $95\%$ CIs within $[0.76, 0.90]$), against $-0.13$ for DINOv3. The self-supervised, low-learning-rate version recovers less than the supervised, high-learning-rate one (mean per-image difference $-0.13$, $95\%$ CI $[-0.22, -0.05]$, Holm $p = 0.005$), and its recoverable component stays in the encoder. Pretraining scale remains confounded: DINO-ResNet-50 is pretrained on ImageNet-1k ($1.28$M images) and DINOv3 on LVD-1689M (Appendix~\ref{app:armL-expA}).

\textbf{Implication.} The recoverable component is architecture-dependent: decoder for DINOv3, encoder for OGLANet, both for MAMNet. This directly explains the catastrophic failures in Table~\ref{tab:loco-main}. MRFP$^+$ perturbs encoder features uniformly without targeting OGLANet's encoder-locus failure, so it costs OGLANet $-14.0$ mIoU. ISW whitens features in a way that corrupts the representation DINOv3 depends on, so it costs DINOv3 $-18.5$ mIoU.

\subsection{External benchmark: S-EO}
\label{sec:seo}

We repeat the benchmark and the three probes on S-EO~\cite{seo}: WorldView-3 imagery of Omaha, Jacksonville and San Diego, four architectures, three seeds, and $3{,}500$ training patches per side (Appendix~\ref{app:seo}). Its masks are derived from digital surface models; both sides of every gap share them, so label noise largely cancels in the gap. The upper bound exceeds LOCO in all $12$ (architecture, held-out city) pairs (Holm $p < 0.05$), by $2.10$ to $7.95$ mIoU and $5.41$ $[5.07, 5.75]$ pooled. Histogram matching lowers mIoU in all $36$ vanilla cells, each significant after Holm correction, and the foreground AURC gap exceeds the background gap in all $12$ pairs, significantly in $8$. Decoder retraining raises foreground IoU in all $18$ MAMNet and DINOv3 cells and lowers it in all $9$ OGLANet cells; the size of this change tracks the gap of the LOCO model whose decoder is retrained ($R^2 = 0.62$), while its sign separates by architecture even between DINOv3 and OGLANet models with similar gaps. The background probability also shifts under LOCO in $10$ of $12$ pairs, so on S-EO the confidence change is not confined to the shadow class. The gap shows no dependence on the sun-elevation distance between test patches and training scenes (cluster-robust $p = 0.30$, three cities).

